\documentclass[twoside]{article}

\usepackage[preprint]{aistats2027}
\usepackage[round,authoryear]{natbib}
\usepackage{amsmath,amssymb,amsthm}
\usepackage{mathtools}
\usepackage{microtype}
\usepackage{booktabs}
\usepackage{balance}
\usepackage{url}

\renewcommand{\bibname}{References}
\renewcommand{\bibsection}{\subsubsection*{\bibname}}
\bibliographystyle{apalike}

\newtheorem{theorem}{Theorem}
\newtheorem{lemma}[theorem]{Lemma}
\newtheorem{corollary}[theorem]{Corollary}
\newtheorem{proposition}[theorem]{Proposition}
\newtheorem{remark}[theorem]{Remark}
\newcommand{\R}{\mathbb R}
\newcommand{\E}{\mathbb E}
\newcommand{\Var}{\operatorname{Var}}
\newcommand{\conv}{\operatorname{conv}}
\newcommand{\supp}{\operatorname{supp}}

\begin{document}
\runningauthor{}

\twocolumn[
\aistatstitle{Sharp Pair Selection for Mean Regression}

\aistatsauthor{Pahan Dewasurendra}

\aistatsaddress{Johns Hopkins University}
]

\begin{abstract}
Suppose a learner may train the empirical mean on two optimally selected
examples, with repetition, and is evaluated by squared loss on the full
dataset. How much worse can this be than the full mean? This is the first open
column in a recent data-selection problem. We determine the exact answer in
every dimension:
\[
 F(d,2)=1+\max\left\{\frac13,\frac{d-1}{2d}\right\}.
\]
Thus the sharp factor is $4/3$ for $d\leq3$ and
$(3d-1)/(2d)$ for $d\geq4$. The two branches reveal distinct obstructions.
A rare point on a line forces the dimension-free term, while a uniform regular
simplex forces the dimension term.

The proof gives a stronger convex-geometric theorem for distributions on at
most $m$ atoms. Its key step is an explicit antithetic distribution over
pairs. If one atom has mass at least $1/4$, a star coupling has exactly one
third of the original second moment. Otherwise, a complete-graph coupling has
a closed-form second moment, and a single Jensen inequality proves that
uniform weights are worst. A variance-minimizing Carath\'eodory reduction then
transfers the result to arbitrary finite datasets. This resolves all cases
with selection budget two, supplies a short proof of the previously isolated
planar case, and separates variance-normalized data selection from
radius-normalized approximate Carath\'eodory bounds.
\end{abstract}

\section{Introduction}

Data selection asks what a fixed learning rule can achieve when an informed
teacher chooses a small training set from a larger population. This viewpoint
appears in machine teaching, sample compression, and coreset construction
\citep{dasgupta2019teaching,feldman2020coresets,hanneke2025selection}. Even
the most elementary regression rule leaves a sharp geometric question open.

Let $D=\{z_1,\ldots,z_N\}$ be a finite multiset in $\R^d$. For squared loss,
write
\[
 L_D(h)=\frac1N\sum_{i=1}^N\|h-z_i\|^2,
 \qquad L_D^*=\min_h L_D(h).
\]
The minimizer is the full mean $\mu_D=N^{-1}\sum_i z_i$. If the rule receives
only $n$ examples, it returns their equal-weight average. Repetitions are
allowed, as in the original formulation. Define
\begin{align*}
 L_D^*(n)&=\min_{x_1,\ldots,x_n\in D}
 L_D\left(\frac1n\sum_{j=1}^n x_j\right),\\
 F(d,n)&=\sup_{D\subset\R^d}\frac{L_D^*(n)}{L_D^*}.
\end{align*}
The ratio is set to one for a constant dataset.

\citet{hanneke2025selection} proved
\[
 F(1,n)=\frac{2n}{2n-1},
 \qquad
 \frac{2n}{2n-1}\leq F(d,n)\leq\frac{n+1}{n},
\]
and showed that the upper bound is approached as $d$ grows. Their COLT open
problem records only one additional finite-dimensional value,
$F(2,2)=4/3$, obtained by a computational geometric argument
\citep{hanneke2025open}. Every other entry with $d,n>1$ was left open.

We close the full $n=2$ column.

\begin{theorem}[Exact pair-selection factor]\label{thm:main}
For every integer $d\geq1$,
\begin{equation}\label{eq:main}
 F(d,2)=\begin{cases}
 4/3,&d\leq3,\\
 (3d-1)/(2d),&d\geq4.
 \end{cases}
\end{equation}
\end{theorem}

For reference, the first values are
\begin{center}
\begin{tabular}{c@{\quad}ccccc@{\quad}c}
\toprule
$d$ & 1 & 2 & 3 & 4 & 5 & $d\to\infty$\\
\midrule
$F(d,2)$ & $4/3$ & $4/3$ & $4/3$ & $11/8$ & $7/5$ & $3/2$\\
\bottomrule
\end{tabular}
\end{center}

The formula is the maximum of two qualitatively different lower bounds. Four
points on a line, three at zero and one at one, force excess risk $1/3$.
The uniform distribution on a regular $d$-simplex forces excess risk
$(d-1)/(2d)$. Their crossover at $d=3$ explains why extrapolating from the
known planar answer misses the higher-dimensional behavior.

Our main technical result is stronger than Theorem~\ref{thm:main}. It gives
the exact constant as a function of the number of atoms, independent of the
ambient Hilbert space.

\begin{theorem}[Sharp atomic theorem]\label{thm:atomic}
Let $m\geq2$. Let $X$ be a finitely supported random vector in a real Hilbert space, with
$\E X=0$, $0<\E\|X\|^2<\infty$, and $|\supp(X)|\leq m$. Then there are
$x,y\in\supp(X)$, possibly equal, such that
\begin{align}\label{eq:atomic}
 \left\|\frac{x+y}{2}\right\|^2
 &\leq C_m\E\|X\|^2,\\
 C_m&=\max\left\{\frac13,\frac{m-2}{2(m-1)}\right\}.
\end{align}
The constant is optimal for every $m\geq2$.
\end{theorem}

The proof is constructive at the level of a randomized certificate. It does
not sample two independent observations. Instead it builds an antithetic pair
whose average has the sharp second moment. This distinction is essential.
Independent sampling gives only $1/2$, which corresponds to the previously
known dimension-free bound $F(d,2)\leq3/2$.

\section{Risk and Atomic Reduction}\label{sec:reduction}

The squared-loss identity
\begin{equation}\label{eq:biasvariance}
 L_D(h)=L_D^*+\|h-\mu_D\|^2
\end{equation}
turns the problem into equal-weight mean approximation. After translating by
$\mu_D$ and normalizing $L_D^*=1$, it is enough to find two data points whose
midpoint has squared norm at most the claimed excess factor.

The following standard reduction makes support size, rather than dataset
size, the relevant parameter. We include the argument because preserving the
variance direction is important here.

\begin{lemma}[Variance-minimizing reduction]\label{lem:reduction}
Let $D\subset\R^d$ be a finite multiset with mean $\mu$ and variance $V$.
There is a probability distribution $P$ supported on at most $r+1$ points of
$D$, where $r$ is the affine dimension of $D$, such that
\[
 \E_P X=\mu,
 \qquad
 \E_P\|X-\mu\|^2\leq V.
\]
\end{lemma}

\paragraph{Proof.}
Minimize $\sum_i p_i\|z_i-\mu\|^2$ subject to $p_i\geq0$,
$\sum_i p_i=1$, and $\sum_i p_i z_i=\mu$. The uniform weights are feasible.
An extreme optimal solution has at most $r+1$ positive coordinates because
there are $r+1$ independent affine constraints. \hfill$\square$

Any pair selected from the reduced support is also available in the original
dataset. Therefore, a bound relative to the reduced variance is at least as
strong as the same bound relative to $V$. Lemma~\ref{lem:reduction} and
Theorem~\ref{thm:atomic} immediately give the upper bound in
Theorem~\ref{thm:main}, since $m\leq d+1$ and $C_m$ is nondecreasing in $m$.

\subsection{A finite minimax formulation}

The reduction exposes a small semidefinite game. Fix positive masses
$p\in\R^m$ and let $\mathcal Q_2$ contain the empirical frequency vectors
$e_i$ and $(e_i+e_j)/2$. Put $D_p=\operatorname{diag}(p)$,
$s=\sqrt p$, and
\[
 u_q=D_p^{-1/2}(q-p)\in s^\perp.
\]
Every centered geometry with unit variance has a Gram matrix $G$ satisfying
$G\succeq0$, $Gp=0$, and $\operatorname{tr}(D_pG)=1$. Equivalently,
$B=D_p^{1/2}GD_p^{1/2}$ is positive semidefinite on $s^\perp$ and has trace
one. Its squared mean error at $q$ is $u_q^\top B u_q$.

\begin{proposition}[Pair minimax duality]\label{prop:minimax}
For a fixed positive mass vector $p$, the worst normalized pair error over all
Hilbert-space configurations is
\begin{align}\label{eq:minimax}
 R(p)
 &=\max_{\substack{B\succeq0,\ Bs=0\\\operatorname{tr}B=1}}
     \min_{q\in\mathcal Q_2}u_q^\top Bu_q\\
 &=\min_{\lambda\in\Delta(\mathcal Q_2)}
   \lambda_{\max}\left(
   \left.\E_{q\sim\lambda}[u_qu_q^\top]\right|_{s^\perp}
   \right).
\end{align}
\end{proposition}

\paragraph{Proof.}
Replace the finite minimum by a minimum over distributions $\lambda$ on
$\mathcal Q_2$. The objective is bilinear in $B$ and $\lambda$, so finite
dimensional minimax duality swaps max and min. Maximizing
$\operatorname{tr}(BM)$ over positive semidefinite trace-one matrices on
$s^\perp$ returns the largest eigenvalue of $M$ on that subspace.
\hfill$\square$

Independent sampling from $p$ makes the second moment in
\eqref{eq:minimax} equal to $(1/2)I$ on $s^\perp$. This recovers excess
$1/2$. The proof below constructs a dependent distribution $\lambda$ with a
strictly smaller spectral second moment. The construction depends only on
$p$, not on the geometry $G$.

\section{The Antithetic Pair Certificate}\label{sec:certificate}

We prove Theorem~\ref{thm:atomic}. Write the support as
$x_1,\ldots,x_m$, with positive masses $p_1,\ldots,p_m$. Center and scale so
that
\begin{equation}\label{eq:normalization}
 \sum_i p_i x_i=0,
 \qquad
 \sum_i p_i\|x_i\|^2=1.
\end{equation}
It is enough to construct a distribution on pairs whose expected squared
midpoint is at most $C_m$. Some pair then attains the same bound.

\subsection{A dominant atom}

Suppose $p_j\geq1/4$ for some $j$. Consider the following distribution on
unordered pairs:
\begin{align}\label{eq:star}
 \Pr\{(j,j)\}&=\frac{4p_j-1}{3},\\
 \Pr\{(i,j)\}&=\frac{4p_i}{3},\qquad i\ne j.
\end{align}
The probabilities are nonnegative and sum to one. If $M$ denotes the midpoint
of the sampled pair, then
\begin{align*}
 3\E\|M\|^2
 &=(4p_j-1)\|x_j\|^2
   +\sum_{i\ne j}p_i\|x_i+x_j\|^2\\
 &=\sum_i p_i\|x_i\|^2=1.
\end{align*}
The second equality uses $\sum_{i\ne j}p_i x_i=-p_jx_j$. Thus a dominant
atom always yields a pair with squared midpoint norm at most $1/3$. The
threshold $1/4$ is exact because it is precisely where the loop probability
in \eqref{eq:star} becomes nonnegative.

\subsection{Diffuse masses}

Now assume $p_i<1/4$ for every $i$. Necessarily $m\geq5$. Define
\begin{align}\label{eq:SCalpha}
 S=\sum_{i=1}^m\frac{p_i}{1-4p_i},
 \qquad C(p)&=\frac{S+1}{3S+1},\\
 a_i&=\frac{C(p)p_i}{(1+S)(1-4p_i)}.
\end{align}
For each distinct pair, set
\begin{equation}\label{eq:complete-coupling}
 w_{ij}=4(p_i a_j+p_j a_i),\qquad i<j.
\end{equation}
All weights are positive. They also sum to one. To see this, let
$A=\sum_i a_i$ and $U=\sum_i p_i a_i$. The identities
\[
 A=\frac{C(p)S}{1+S},
 \qquad
 U=\frac{C(p)(S-1)}{4(1+S)}
\]
give $\sum_{i<j}w_{ij}=4(A-U)=1$.

\begin{lemma}[Exact diffuse certificate]\label{lem:diffuse}
If a distinct pair $(i,j)$ is drawn with probabilities
\eqref{eq:complete-coupling}, then its midpoint $M$ satisfies
\begin{equation}\label{eq:exact-certificate}
 \E\|M\|^2=C(p)\sum_i p_i\|x_i\|^2=C(p).
\end{equation}
\end{lemma}

\paragraph{Proof.}
Let $q=(e_i+e_j)/2$ encode the sampled midpoint, so $M=\sum_i q_i x_i$.
Write $Q=\E[qq^\top]$. The off-diagonal entries from
\eqref{eq:complete-coupling} are
\[
 Q_{ij}=p_i a_j+a_i p_j.
\]
The degree identity
\[
 \frac14\sum_{j\ne i}w_{ij}=C(p)p_i+2p_i a_i
\]
follows from
$a_i(1-4p_i)=p_i(C(p)-A)$. Hence
\begin{equation}\label{eq:Qidentity}
 Q=C(p)\operatorname{diag}(p)+pa^\top+ap^\top.
\end{equation}
If $G_{ij}=\langle x_i,x_j\rangle$, centering gives $Gp=0$. Therefore
\[
 \E\|M\|^2=\operatorname{tr}(GQ)
 =C(p)\operatorname{tr}(G\operatorname{diag}(p))=C(p).
\]
\hfill$\square$

It remains to maximize $C(p)$. The function $t\mapsto t/(1-4t)$ is strictly
convex on $(0,1/4)$. Jensen's inequality gives
\begin{equation}\label{eq:jensen}
 S\geq\frac{m}{m-4}.
\end{equation}
Since $(S+1)/(3S+1)$ decreases with $S$,
\begin{equation}\label{eq:Cbound}
 C(p)\leq
 \frac{m/(m-4)+1}{3m/(m-4)+1}
 =\frac{m-2}{2(m-1)}.
\end{equation}
Equality in \eqref{eq:jensen} holds exactly at $p_i=1/m$. Combining the
dominant and diffuse cases proves the upper bound in
Theorem~\ref{thm:atomic}.

The diffuse certificate is not merely an upper bound. It solves the fixed
mass-profile minimax problem exactly.

\begin{theorem}[Exact diffuse mass profile]\label{thm:profile}
Fix positive masses $p_1,\ldots,p_m<1/4$ that sum to one. Over all centered
Hilbert-space configurations with unit variance,
\begin{equation}\label{eq:profile}
 R(p)=C(p)=
 \frac{1+\sum_i p_i/(1-4p_i)}
 {1+3\sum_i p_i/(1-4p_i)}.
\end{equation}
The supremum is attained in dimension $m-1$ by an explicit simplex.
\end{theorem}

\paragraph{Proof.}
Lemma~\ref{lem:diffuse} proves $R(p)\leq C(p)$. For the reverse inequality,
write $C=C(p)$ and define
\begin{equation}\label{eq:profile-diagonal}
 b_i=C+\frac{3C-1}{1-4p_i}.
\end{equation}
Let $G$ have diagonal $G_{ii}=b_i$ and, for $i\ne j$,
\begin{equation}\label{eq:profile-gram}
 G_{ij}=2C-\frac{b_i+b_j}{2}.
\end{equation}
First, $\sum_i p_i b_i=1$. This follows from
$3C-1=2/(3S+1)$ and the definition of $C$. The identity
\[
 (1-4p_i)b_i=4C(1-p_i)-1
\]
then gives $Gp=0$ by a row-wise calculation.

The matrix is positive semidefinite. Indeed, for every $y$ with
$\sum_i y_i=0$, Equations~\eqref{eq:profile-diagonal} and
\eqref{eq:profile-gram} give
\begin{equation}\label{eq:profile-psd}
 y^\top Gy=2\sum_i(b_i-C)y_i^2>0.
\end{equation}
For an arbitrary $z$, set $y=z-(\sum_i z_i)p$. Then $\sum_i y_i=0$, and
$z^\top Gz=y^\top Gy$ because $Gp=0$. Hence $G\succeq0$, with rank $m-1$.
It is therefore the Gram matrix of vectors $x_1,\ldots,x_m$ with weighted
mean zero and weighted variance one.

Every repeated pair has cost $b_i\geq C$. By
\eqref{eq:profile-gram}, every distinct pair has cost exactly
\[
 \frac{G_{ii}+G_{jj}+2G_{ij}}4=C.
\]
Thus the best pair in this configuration has cost $C$, proving
\eqref{eq:profile}. \hfill$\square$

Theorem~\ref{thm:profile} describes a continuum of least favorable
nonregular simplices. As a mass approaches $1/4$ from below, its squared
vertex norm approaches $3$ under our normalization and $C(p)$ approaches
$1/3$. At uniform masses, the Gram matrix becomes the regular-simplex matrix
and $C(p)$ reaches the diffuse maximum in \eqref{eq:Cbound}.

\section{Sharpness and Extremizers}\label{sec:sharpness}

The lower bounds are elementary and expose why both branches are necessary.

\paragraph{Rare-point obstruction.}
Let $X$ equal zero with probability $3/4$ and one with probability $1/4$.
Its mean is $1/4$ and its variance is $3/16$. Every average of two support
points lies in $\{0,1/2,1\}$. Its squared distance from the mean is at least
$1/16$, giving normalized excess $1/3$. The finite multiset
$\{0,0,0,1\}$ realizes this example exactly.

\paragraph{Simplex obstruction.}
Let $x_1,\ldots,x_m$ be the unit vertices of a regular $(m-1)$-simplex,
so $\sum_i x_i=0$ and
$\langle x_i,x_j\rangle=-1/(m-1)$ for $i\ne j$. Under uniform masses the
variance is one. Repeating a vertex gives squared norm one, while every
distinct midpoint has squared norm
\[
 \frac14\left(2-\frac{2}{m-1}\right)
 =\frac{m-2}{2(m-1)}.
\]
This proves sharpness of Theorem~\ref{thm:atomic}. With $m=d+1$, it gives the
second lower bound in Theorem~\ref{thm:main}.

The diffuse proof also characterizes the high-dimensional extremizer.

\begin{proposition}[Uniqueness on a minimal support]\label{prop:unique}
Fix $m\geq5$. If a centered distribution on exactly $m$ atoms attains the
constant $(m-2)/(2(m-1))$ in Theorem~\ref{thm:atomic}, then its masses are
uniform and its support is a regular $(m-1)$-simplex, up to scaling and an
isometry on its span.
\end{proposition}

\paragraph{Proof.}
The dominant certificate is strictly smaller than the claimed value. In the
diffuse case, equality in \eqref{eq:jensen} forces uniform masses. The weights
$w_{ij}$ then become uniform over distinct pairs. Since their average cost
equals the minimum, all distinct midpoint costs are equal. Normalize the
variance to one and write $b_i=\|x_i\|^2$. Centering, equal midpoint costs,
and $\sum_i b_i=m$ imply
\[
 0=\left\langle x_i,\sum_jx_j\right\rangle
 =\frac{m-4}{2}(1-b_i).
\]
Thus every $b_i=1$, after which all off-diagonal inner products equal
$-1/(m-1)$. \hfill$\square$

\section{Consequences}\label{sec:consequences}

The proof provides more information than the worst-case constant. For any
variance-minimizing representation with masses $p$, it gives an explicit
certificate. A mass at least $1/4$ guarantees excess $1/3$. If all masses are
smaller, the sharper data-dependent value is
\begin{equation}\label{eq:adaptive}
 C(p)=\frac{1+\sum_i p_i/(1-4p_i)}
 {1+3\sum_i p_i/(1-4p_i)}.
\end{equation}
This interpolates continuously between the dominant-atom boundary and the
uniform-simplex maximum. It also supplies a direct witness distribution over
good pairs, rather than only a minimax value.

Lemma~\ref{lem:reduction} gives an intrinsic-dimension version. If a dataset
has affine dimension $r$, regardless of its ambient dimension, two examples
achieve
\begin{equation}\label{eq:intrinsic}
 \frac{L_D^*(2)}{L_D^*}
 \leq1+\max\left\{\frac13,\frac{r-1}{2r}\right\}.
\end{equation}
The bound is sharp for every $r$. The exact factor approaches $3/2$ from
below, quantifying the finite-dimensional improvement over independent
sampling.

The certificate has a simple semidefinite interpretation. For masses $p$,
let $\Sigma_p=\operatorname{diag}(p)-pp^\top$. A distribution on empirical
two-point frequencies $q$ proves a uniform geometric bound $C$ whenever
\[
 \E[(q-p)(q-p)^\top]\preceq C\Sigma_p.
\]
The star construction attains equality with $C=1/3$. In the diffuse case,
Equation~\eqref{eq:Qidentity} is the same certificate after quotienting out
the centered direction $p$. This explains why the pair distribution can be
chosen without knowing the geometry of the support.

The exact profile also gives a stability statement for the diffuse branch.
For $m\geq5$, strict convexity in \eqref{eq:jensen} implies that any sequence
of diffuse mass vectors whose minimax values approach
$(m-2)/(2(m-1))$ must converge, up to permutation, to the uniform vector.
Together with Proposition~\ref{prop:unique}, near-worst minimal-support
instances must approach the regular simplex in both masses and Gram geometry.

\section{Relation to Approximate Carath\'eodory Theory}

Approximate Carath\'eodory theorems replace a convex combination by a sparse
or equal-weight combination. Maurey's empirical method gives the familiar
$O(1/\sqrt n)$ Euclidean error by independent sampling. Later work gives
constructive algorithms, tight dimension-free rates, and discrepancy-based
extensions \citep{mirrokni2017tight,combettes2023revisiting,
reis2022discrepancy}.

The normalization is decisive. Those results control error by a containing
radius or diameter. Our denominator is the variance under the particular
barycentric weights, which can be much smaller than the squared radius when a
distant vertex has low mass. The rare-point obstruction in
Section~\ref{sec:sharpness} is exactly this gap.

The closest result is the sharp dimension-dependent theorem of
\citet{tinarrage2026sharp}. For a simplex in the unit ball, the worst squared
distance to a two-fold barycenter is $(d-1)/(2d)$ for $d\geq2$, attained by a
regular simplex. That theorem settles the radius-normalized problem through a
Gram-matrix optimization. It does not imply a variance-normalized bound. Our
formula shows precisely what changes after variance normalization. The same
regular simplex remains extremal for $d\geq4$, while a one-dimensional rare
point dominates in dimensions at most three.

Regular simplices also uniquely maximize variance among distributions with a
fixed support diameter \citep{lim2022variance}. That isodiametric result varies
the distribution and support under a diameter constraint. Here the variance
is fixed, pair midpoints are the admissible decisions, and the extremizer
changes with dimension.

Regular-grid rounding on the standard simplex has also been solved for broad
classes of coordinate norms \citep{bomze2014rounding}. In our formulation,
the norm on barycentric error is induced by the unknown support Gram matrix,
while the normalization is induced by its weighted trace. Proposition
\ref{prop:minimax} makes this adversarial metric explicit.

\section{Discussion}

Theorem~\ref{thm:main} resolves one complete column of the data-selection
table and gives a short analytic proof of its previously known planar entry.
The proof identifies a dominant-or-diffuse dichotomy. Dominant atoms are
handled by a star coupling. Diffuse masses are handled by a complete-graph
coupling, whose worst case follows from scalar convexity.

For budgets larger than two, the admissible averages form a finer simplex
grid. The same second-moment viewpoint leads to distributions over count
vectors, but pair weights are replaced by higher-order inclusion
probabilities. Determining their sharp variance-normalized constants remains
open. The present result suggests that any full formula must retain at least
two competing families: rare lower-dimensional distributions and balanced
regular simplices.

\balance
\bibliography{references}

\appendix
\newpage
\section{Pointers to the Supplement}
The supplement gives a standalone proof of the atomic theorem, a matrix form
of both antithetic certificates, the intrinsic-dimension reduction, and
additional details on sharpness and equality.

\end{document}
